\par\noindent\textbf{(1)} 基取 \((|2,0\rangle,|1,1\rangle,|0,2\rangle)\)。相互作用能分别为 \(U,0,U\)，跃迁矩阵元由玻色升降算符给 \(-\sqrt2J\) 在相邻Fock态之间，故矩阵 \(H=\begin{pmatrix}U&-\sqrt2J&0\\-\sqrt2J&0&-\sqrt2J\\0&-\sqrt2J&U\end{pmatrix}\)。
\par\noindent\textbf{(2)} 反对称态 \(|A\rangle=(|2,0\rangle-|0,2\rangle)/\sqrt2\) 与 \(|1,1\rangle\) 解耦，能量 \(E_A=U\)。对称态 \(|S\rangle=(|2,0\rangle+|0,2\rangle)/\sqrt2\) 与 \(|1,1\rangle\) 的块为 \(\begin{pmatrix}U&-2J\\-2J&0\end{pmatrix}\)。
\par\noindent\textbf{(3)} 两个本征能 \(E_\pm=[U\pm\sqrt{U^2+16J^2}]/2\)。当 \(U\gg J\) 时基态 \(E_-\approx-4J^2/U\)，其主要为单站点各一个粒子的 \(|1,1\rangle\)，双占据仅为 \(O[(J/U)^2]\)，体现相互作用阻塞。
\par\noindent\textbf{(4)} 当 \(U=0\)，对称块本征能为 \(\pm2J\)，反对称态为零。任何初始 \(|2,0\rangle\) 既含对称又含反对称分量，动力学具有能差 \(2J/\hbar\) 与 \(4J/\hbar\) 的相干成分；若加入随机相位或退相干，振荡可衰减但并非此封闭Hamilton模型本身造成。
